\section{Concluding statements}\label{sec:conclusion52}
Within the stated orthogonal numerical interface, a fixed positive
error below $\rho/(2D)$ does not permit uniformly bounded hidden width
for an infinite generated group. The proof combines finite-orbit
reduction with a wordwise-calibrated contraction and allows arbitrary
hidden observation fibers. For rational rotations it yields an
explicit logarithmic lower rate. The independent section geometry now
also has an effective one-surplus antipodal multiplier and a numerical
sufficient six-state horizon. Finite algebraic separation propagates
one positive noise interval to all larger horizons, while reachable-basis
conditioning gives a separate explicit interval.

The remaining distinctions are precise. The universal noise boundary
is not sharp; the intermediate region between the proved small-error
range and the trivial fair-answer range is not classified. The
rational-height rate need not be optimal, and no sharp general formula
for $\gamma_{\calA}(V,\rho)$ is established. The unconditional
six-state noise constant is computable but has not been numerically
evaluated at the large illustrative horizon. Higher-surplus sharp
frontiers, partial-query interfaces, nonspanning seeds and
nonorthogonal command families require their own statements. The
complete global certificate hierarchy is not a practical general
optimizer merely because it converges.

These are boundaries of the conclusions, not alterations to their
hypotheses or substitutions for earlier results. The retained appendix
contains the earlier proofs in full. The finite-dimensional dependency
chain developed here does not establish the independent analytic
Fourier, stopped large-deviation, filtering or nonlinear-resolvent
gates of the wider General Theta series.
